\documentclass[a4paper,12pt]{article}

% Pacchetti per la lingua e la codifica
\usepackage[utf8]{inputenc}
\usepackage[T1]{fontenc}
\usepackage[italian]{babel}

% Pacchetti per la formattazione e il layout
\usepackage{geometry}
\geometry{a4paper, margin=2.5cm}
\usepackage{microtype}
\usepackage{parskip} % Spaziatura tra i paragrafi
\usepackage{setspace}
\onehalfspacing

% Pacchetti per elementi grafici e matematici
\usepackage{amsmath, amssymb}
\usepackage{graphicx}
\usepackage{xcolor}
\usepackage{tcolorbox}
\usepackage{hyperref}
\usepackage{enumitem}

% Configurazione dei link
\hypersetup{
    colorlinks=true,
    linkcolor=blue,
    filecolor=magenta,      
    urlcolor=cyan,
    pdftitle={Appunti Cybersecurity - Introduzione},
    pdfpagemode=FullScreen,
}

% Definizione di un box per le definizioni importanti
\newtcolorbox{definizione}[1]{
  colback=blue!5!white,
  colframe=blue!75!black,
  title={\textbf{#1}}
}

% Definizione di un box per le formule (come richiesto)
\newtcolorbox{formula}{
  colback=white,
  colframe=black,
  sharp corners,
  boxrule=1pt
}

% Titolo del documento
\title{\textbf{Introduzione alla Cybersecurity}\\
\large Appunti del corso basati sulle slide ufficiali}
\author{Prof. Francesco Palmieri}
\date{}

\begin{document}

\maketitle
\tableofcontents
\newpage

\section{Panoramica del Corso}

Lo scopo principale del corso è costruire un ecosistema di competenze adeguate per gestire strumenti e architetture a supporto della sicurezza informatica. Verranno chiariti i concetti fondamentali, analizzate le potenziali minacce e studiate le tecnologie disponibili per fronteggiarle.

Il corso si articola in:
\begin{itemize}
    \item \textbf{Lezioni frontali:} 7 CFU (56 ore)
    \item \textbf{Laboratorio:} 2 CFU (16 ore)
\end{itemize}

L'esame consiste in un colloquio individuale. I prerequisiti fondamentali includono i fondamenti di sicurezza, reti di calcolatori e familiarità con i sistemi Unix/Linux.

\section{Etica e Legalità}

Durante il corso verranno discussi e simulati diversi tipi di attacchi, utilizzando tecnologie di intercettazione e scansione delle vulnerabilità. È fondamentale comprendere il confine etico e legale di tali attività.

\begin{definizione}{Regola Fondamentale}
Nessuna tecnica o strumento appreso deve essere utilizzato su reti reali senza il consenso informato di tutte le parti coinvolte. L'esistenza di vulnerabilità non giustifica accessi non autorizzati.
\end{definizione}

L'obiettivo è formare professionisti (\textit{White Hat}) in grado di difendere le infrastrutture, e non attaccanti (\textit{Black Hat}) che ne violano l'integrità. Il mancato rispetto delle norme può portare a gravi conseguenze legali.

\section{Definizioni Fondamentali}

\subsection{Cyberspazio e Cybersecurity}

Il \textbf{cyberspazio} è definito come l'insieme dei sistemi informatici e delle reti di interconnessione. È un'opera estremamente complessa, costituita da milioni di reti, stratificazioni di software e dispositivi eterogenei. Questa complessità intrinseca genera vulnerabilità che possono essere sfruttate per lanciare attacchi, con conseguenze sia nel mondo digitale che in quello reale.

\begin{definizione}{Cybersecurity}
La Cybersecurity è la sicurezza dei sistemi informatici e delle reti di interconnessione, con lo scopo di proteggerne l'operatività e gli asset strategici in caso di attacchi, incidenti o guasti, garantendo disponibilità, affidabilità, integrità e riservatezza.
\end{definizione}

\subsection{I tre pilastri della Cybersecurity (CIA Triad)}

La sicurezza informatica si fonda su tre concetti chiave:

\begin{enumerate}
    \item \textbf{Riservatezza (Confidenzialità):} Garantire che informazioni e risorse siano accessibili solo a soggetti autorizzati. Riguarda la protezione dei contenuti, la privacy degli individui e la segretezza delle organizzazioni.
    \item \textbf{Integrità:} Garantire che le informazioni non siano state modificate, manipolate o distrutte in modo non autorizzato. Si distingue in:
    \begin{itemize}
        \item \textit{Integrità dei dati:} le informazioni vengono modificate solo in modo autorizzato.
        \item \textit{Integrità del sistema:} il sistema opera libero da manipolazioni esterne.
    \end{itemize}
    \item \textbf{Disponibilità:} Assicurare che gli utenti legittimi abbiano sempre accesso alle risorse quando necessario. 
\end{enumerate}

Matematicamente, per la disponibilità si desidera che:

\begin{formula}
\[
\lim_{t \to \infty} D(t) = 1
\]
\end{formula}
\textit{Dove $D(t)$ è la probabilità che il sistema funzioni correttamente all'istante $t$.}

\subsection{Concetti correlati}

Oltre ai tre pilastri, vi sono altre proprietà essenziali:
\begin{itemize}
    \item \textbf{Autenticità:} Certezza che un dato provenga da una specifica sorgente.
    \item \textbf{Non ripudiabilità:} Impossibilità per un soggetto di negare di aver compiuto una determinata azione (es. invio di un messaggio).
    \item \textbf{Resilienza:} Capacità di un sistema di operare anche in condizioni degradate e di recuperare l'operatività in tempi accettabili dopo un incidente.
\end{itemize}

\section{Safety, Security e Dependability}

È importante distinguere tra concetti spesso confusi:

\begin{itemize}
    \item \textbf{Safety:} Capacità di un sistema di funzionare senza provocare danni alle persone o all'ambiente (es. sensori che aprono le porte di un'auto in caso di incidente).
    \item \textbf{Security:} Protezione delle risorse informatiche da accessi non autorizzati o danni intenzionali (es. evitare che qualcuno apra le porte dell'auto saltando sul tetto, sfruttando il sensore di safety).
    \item \textbf{Dependability (Affidabilità):} Proprietà che permette di fare assoluto affidamento sul servizio fornito.
\end{itemize}

La \textbf{Dependability} si analizza attraverso tre componenti:
\begin{enumerate}
    \item \textbf{Minacce:} Circostanze che portano alla perdita di affidabilità.
    \item \textbf{Attributi:} Proprietà per valutare la qualità del servizio.
    \item \textbf{Strumenti:} Metodi tecnici per garantire il servizio.
\end{enumerate}

\section{Il contesto Internet e le sue debolezze}

Internet è diventata un'infrastruttura critica per la società (e-commerce, e-health, servizi pubblici), ma il suo modello di sicurezza è rimasto quello degli anni '70. 
Le principali debolezze strutturali includono:
\begin{itemize}
    \item Assenza di un'autorità centrale di controllo.
    \item Impossibilità di un tracciamento sistematico delle connessioni.
    \item Perimetri di rete "porosi" o assenti ("borderless"), specialmente con il wireless.
    \item Tecnologie non nate per transazioni \textit{mission-critical}.
\end{itemize}

Inoltre, le \textbf{bad practices} degli utenti (password deboli, disattivazione antivirus, phishing) rendono la sicurezza vulnerabile nel suo "anello più debole". L'approccio della "sicurezza personale" (proteggo solo il mio PC) è fallimentare: la sicurezza in rete richiede cooperazione e interscambio di informazioni.

\section{Nuovi Scenari di Minaccia: IoT e Cyber-Fisico}

L'espansione della rete ha portato all'\textbf{Internet of Things (IoT)}. Si stima che entro pochi anni ci saranno 50 miliardi di dispositivi connessi.

\subsection{Cyber-Physical Systems}
Le minacce non riguardano più solo i dati, ma anche il mondo fisico (Safety). Esempi critici includono:
\begin{itemize}
    \item \textbf{Dispositivi medici:} Pacemaker e defibrillatori vulnerabili (es. caso Abbott 2017).
    \item \textbf{Veicoli e trasporti:} Controllo remoto non autorizzato di automobili o gru industriali.
    \item \textbf{Domotica:} Termostati e videocamere accessibili pubblicamente.
\end{itemize}
Motori di ricerca come \textbf{Shodan} permettono di individuare facilmente dispositivi IoT vulnerabili connessi in rete.

\section{Tipologie e Motivazioni degli Attacchi}

Gli attacchi si sono evoluti da semplici dimostrazioni di abilità a vere e proprie industrie criminali. 

\subsection{Motivazioni}
\begin{itemize}
    \item \textbf{Guadagno economico:} Furto di dati bancari, ransomware, truffe.
    \item \textbf{Spionaggio:} Furto di proprietà intellettuale o segreti di stato.
    \item \textbf{Vendetta o Attivismo (Hacktivism):} Attacchi ideologici.
    \item \textbf{Cyber Warfare:} Conflitti tra nazioni.
\end{itemize}

\subsection{Vettori di attacco comuni}
Il flusso operativo tipico include:
\begin{enumerate}
    \item \textbf{Malware e Phishing:} Per ottenere l'accesso iniziale.
    \item \textbf{Botnet:} Creazione di reti di dispositivi "zombie".
    \item \textbf{DDoS (Distributed Denial of Service):} Per inibire i servizi.
    \item \textbf{Esfiltrazione dati:} Furto di database e credenziali.
\end{enumerate}

\section{Infrastrutture Critiche}

Le infrastrutture critiche (energia, trasporti, finanza) dipendono dai sistemi informatici e sono bersagli primari.

\subsection{Smart Grids (Rete Elettrica)}
Le reti elettriche moderne integrano flussi di energia e informazioni. Utilizzano sistemi SCADA spesso obsoleti e vulnerabili. Un attacco può causare blackout su vasta scala o danni fisici alle apparecchiature.

\subsection{Controllo del Traffico Aereo (ATC)}
Dipende da comunicazioni radio e dati radar. Rischi includono:
\begin{itemize}
    \item \textbf{Transponder Spoofing:} Creazione di aerei "fantasma".
    \item \textbf{Radar Jamming:} Accecamento dei radar.
\end{itemize}

\subsection{GPS (Global Positioning System)}
Fondamentale per la navigazione e il timing delle reti. È vulnerabile a:
\begin{itemize}
    \item \textbf{Jamming:} Disturbo del segnale per impedirne la ricezione.
    \item \textbf{Spoofing:} Invio di falsi segnali per alterare la posizione percepita (facilitato oggi da dispositivi \textit{Software Defined Radio} a basso costo).
\end{itemize}

\section{Cyber Warfare e Terrorismo}

Il cyberspazio è diventato il quinto dominio bellico (dopo terra, mare, aria e spazio). 

\begin{definizione}{Cyber Warfare}
Attacchi informatici contro l'integrità o la disponibilità di sistemi critici, tali da costituire una minaccia per la sicurezza nazionale, spesso a supporto di attacchi fisici (cinetici).
\end{definizione}

Un esempio storico è \textbf{Stuxnet} (2010), un worm progettato per sabotare selettivamente le centrifughe di arricchimento dell'uranio iraniane, modificando la logica dei controllori industriali (PLC) senza causare danni ai computer infettati durante la propagazione.

\section{Conclusioni: La Sfida}

La probabilità di subire un attacco è funzione del numero di utenti, dei confini territoriali, delle problematiche tecniche e della disponibilità di informazioni.

Non esiste una "soluzione definitiva" (silver bullet). La sicurezza è un processo continuo che richiede:
\begin{itemize}
    \item Consapevolezza del problema.
    \item Competenze adeguate.
    \item Cooperazione tra governi, organizzazioni e utenti.
    \item Integrazione di tecnologia, fattori economici, influenza sociale e politiche normative.
\end{itemize}

\end{document}
